\section*{Chapitre \ref{ch:in:the-options-market-making-houses} --- Les maisons de market making d'options}

\begin{solution}{exo:in:the-options-market-making-houses:1}
$15\times40\times2=1\,200$ séries.
\end{solution}

\begin{solution}{exo:in:the-options-market-making-houses:2}
$1\,158/2\,773=41.8~\%$.
\end{solution}

\begin{solution}{exo:in:the-options-market-making-houses:3}
IMC (la bourse d'options sur actions d'Amsterdam), Susquehanna (la Philadelphia Stock Exchange), Belvedere Trading (le
Chicago Board Options Exchange).
\end{solution}

\begin{solution}{exo:in:the-options-market-making-houses:4}
$1\,369/1\,158-1=18.2~\%$ en 2024 et $1\,769/1\,369-1=29.2~\%$ en 2025. En 2024, le revenu net de trading progresse de 26~\%~: les coûts
(salaires, technologie, impôts) augmentent plus vite que le revenu, si bien que la marge recule de 41.8~\% à 39.2~\%. Les communiqués ne détaillent pas
les coûts~; on ne sait donc pas, à partir d'eux, quel coût a augmenté.
\end{solution}

\begin{solution}{exo:in:the-options-market-making-houses:5}
$4\,556/2\,400=1.90$ à $4\,556/2\,000=2.28$~: 1.90 à 2.28 millions d'euros par tête.
\end{solution}

\begin{solution}{exo:in:the-options-market-making-houses:6}
Son livre est une surface de volatilité, de skew et d'expositions à la décroissance temporelle, qui ne peut être couverte
qu'en partie et lentement, en négociant d'autres options~; et ses obligations de cotation la conduisent à prendre des positions qu'elle ne choisirait pas.
\end{solution}

\begin{solution}{exo:in:the-options-market-making-houses:7}
30.4~\% en 2024 et 34.0~\% en 2025~; $4\,556/2\,500=\text{\euro}1.82$ million par tête.
\end{solution}

\begin{solution}{exo:in:the-options-market-making-houses:8}
Une maison qui choisit de publier n'est pas un échantillon de l'industrie~: elle peut figurer parmi les plus grandes et les plus performantes~;
son effectif est une borne inférieure, si bien que son chiffre est une borne supérieure~; et 2024--2025 furent des années de gros volumes
de l'aveu même de la maison. Le revenu par tête varie selon la firme, selon l'année et selon ce qui est compté.
\end{solution}

\begin{solution}{pb:in:the-options-market-making-houses:1}
\begin{enumerate}
\item Une \omterm{def:in:a-map-of-the-industry:ptf}{firme de trading pour compte propre} (principal trading firm) dont l'activité principale consiste à coter des options cotées sur plusieurs strikes et échéances, et
à les couvrir~; le livre~1 présente le teneur de marché d'options comme le rôle tenu sur une plateforme, avec ses obligations.
\item La bourse d'options sur actions d'Amsterdam (IMC), la Philadelphia Stock Exchange (Susquehanna), le Chicago Board Options Exchange
(Belvedere Trading).
\item Années 1980~: 4, années 1990~: 1, années 2000~: 1.
\item Une firme fondée par des personnes qui quittent une firme existante pour y faire la même activité~; aucune firme n'affiche une telle origine sur
ses propres pages.
\item Elle exploite un marché hybride qui combine la \omterm{def:in:proprietary-market-makers-and-high-frequency-firms:floor}{criée} (open outcry) et le trading électronique.
\item $20\times50\times2=2\,000$.
\item Le nombre d'instruments, le risque détenu, les obligations, le modèle (une surface, non un prix).
\item La tarification, les grecques, les surfaces, l'inventaire et la couverture~; la pratique passe par des jeux, des marchés simulés et
un travail encadré.
\item Des centaines de milliers de cotations bougent à chaque changement du sous-jacent~; aucune personne ne peut les mettre à jour.
\item Le contrat de travail.
\item Revenu net de trading de 2\,773, 3\,494 et 4\,556 millions d'euros~; résultat net de 1\,158, 1\,369 et 1\,769 millions d'euros.
\item 26.0~\% et 30.4~\%.
\item 41.8~\%, 39.2~\%, 38.8~\%.
\item 30.4~\% et 34.0~\%.
\item 1.66 million d'euros en 2024~; au plus 2.28 millions d'euros en 2025.
\item 4\,556 millions d'euros, 38.8~\% et 34.0~\%~; trois des six firmes (50~\%) rattachent leur fondation à une \omterm{def:in:proprietary-market-makers-and-high-frequency-firms:floor}{criée}.
\item Plus de 2\,000 salariés, c'est diviser par au moins 2\,000.
\item Elle a choisi de publier, elle ne représente qu'une firme, et ses années furent chargées~; il n'y a pas d'échantillon.
\item Les comptes de sa filiale britannique, ou la déclaration publique de son broker-dealer américain (chapitre~2).
\item Les maisons forment au jugement sur la tarification et le risque sur une surface, qui est un métier différent de la course à un
prix~; les chiffres d'une seule maison montrent que l'activité peut être très rentable, non ce que gagne une maison donnée.
\end{enumerate}
\end{solution}

\begin{solution}{iq:in:the-options-market-making-houses:1}
Son inventaire est un ensemble d'options dont le risque (vega, skew, gamma) ne peut pas être aplati rapidement sans payer la fourchette
sur de nombreuses séries~; elle couvre le delta en continu et conserve le reste, qu'elle gère sur plusieurs jours.
\iqlookfor{la différence entre couvrir le delta et aplatir un livre de grecques.}
\end{solution}

\begin{solution}{iq:in:the-options-market-making-houses:2}
$12\times30\times2=720$ séries, deux côtés chacune~: 1\,440 cotations.
\iqlookfor{calls et puts, cours acheteurs et cours vendeurs.}
\end{solution}

\begin{solution}{iq:in:the-options-market-making-houses:3}
Pour les avantages de la désignation~: remises sur frais, avantages d'allocation ou de priorité, et le flux qui vient avec le fait d'être
présent (livre~1, chapitre~24). Le coût est d'être présent sur des séries et à des moments que la firme éviterait, avec une sélection adverse
et du capital immobilisé.
\iqlookfor{des obligations achetées au prix de privilèges.}
\end{solution}

\begin{solution}{iq:in:the-options-market-making-houses:4}
Environ $1.29/1.12\approx1.15$~: rentabilité des fonds propres en hausse d'environ 15~\%. Il faut vérifier si l'année a été
anormalement volatile, si la croissance vient d'une seule activité ou d'un seul événement, et si le risque pris a augmenté.
\iqlookfor{le calcul des rapports et le cycle avant le talent.}
\end{solution}

\begin{solution}{iq:in:the-options-market-making-houses:5}
Comment se comportent les modèles et les surfaces propres à la firme sur son propre flux~; comment gérer un livre réel sous limites~; comment
fonctionnent les outils et les alertes du desk~; la microstructure du marché sur des plateformes précises.
\iqlookfor{la pratique, les outils et le jugement, non la théorie.}
\end{solution}

\begin{solution}{iq:in:the-options-market-making-houses:6}
La gamma longue rapporte grâce aux mouvements mais se déprécie~; la vega courte perd si la volatilité implicite monte. Il faut vérifier la couverture en delta, le
calendrier des événements de nuit et la vega par échéance~; réduire la vega si un événement peut faire monter la volatilité implicite, et
décider de la gamma à conserver, au regard du mouvement de nuit attendu et de la dépréciation.
\iqlookfor{les grecques comme risques assortis de prix, et un plan pour la nuit.}
\end{solution}
